\documentclass[10pt,a4paper]{amsart}
\usepackage[margin=19mm]{geometry}
\usepackage{amsmath,amssymb,mathtools}
\usepackage[scr=rsfs,cal=euler]{mathalfa}
\usepackage{xfrac}
\usepackage[unicode,hidelinks,bookmarks=false]{hyperref}
\newcommand{\ie}{i.e.\ }
\newcommand{\cf}{cf.\ }
\newcommand{\resp}{respectively\ }
\newcommand{\Subsection}{\subsection}
\providecommand{\nobreakdash}{\nobreak\hbox{-}}
\setlength{\emergencystretch}{2em}
\multlinegap0pt
\usepackage{amssymb}
\usepackage{xcolor}
\usepackage{enumerate}
\usepackage{bm}
\usepackage{dsfont}
\usepackage{mathtools}
\newtheorem{proposition}{Proposition}
\title{Colour algebras over rings}
\author{S. Pumpluen}
\date{Source: arXiv:2409.14574v2 (CC BY 4.0)}
\begin{document}
\maketitle

\noindent\emph{Source and licence.} Colour algebras over rings. Version arXiv:2409.14574v2 (CC BY 4.0), distributed by arXiv under the Creative Commons Attribution 4.0 International licence, http://creativecommons.org/licenses/by/4.0/, as recorded in the arXiv licence metadata for this exact version and shown on the abstract page https://arxiv.org/abs/2409.14574v2. Full licence notice: \texttt{licenses/arxiv-2409.14574-CC-BY-4.0.txt}.\par\smallskip

\noindent\textit{Editorial scope.} Counted unit: the article's proposition proposition (source0.tex lines 532--543). The article's complete original proof of it is delivered with it (source0.tex lines 545--564, one \texttt{proof} environment of 20 lines). No mathematics is written, repaired or re-derived: the statement and the proof are the article's own lines, in the article's own notation, and no retained line is substituted or re-flowed.  No further source-local context is needed: every cross-reference of every retained line resolves inside the two delivered spans.  Nothing was omitted from any retained span, so the package carries no omission record.  No retained line contains a citation, so the wrapper carries no bibliography and no bibliography entry.  No retained line contains a figure environment, an image-inclusion command or drawing code, so no image is delivered and the package declares no asset dependency.  Editorial formatting only, in addition: an AMS \texttt{amsart} preamble that restates the article's own declarations and macros used by the retained lines, namely the environments \texttt{proposition} on separate counters, together with the standard packages those lines need.  The retained environments print in this document as Proposition 1 in order of appearance, whereas the article prints the same items with the numbering of its own full document.  The complete native source of the article as distributed in the arXiv e-print for this exact version is bundled unchanged as \texttt{sources/165636/source0.tex}.
 \begin{proposition}
 Let $ \varphi:(P,\alpha)\rightarrow (P',\alpha') $
 be a morphism that is an isometry between $h$ and $h'$.
\\ (i)  $\varphi(u\times_{\alpha}v)=u\times_{\alpha'}v$.
 \\ (ii)  $W(S,P,h,\alpha)\cong W(S, P',h',\alpha')$
 via $\varphi$.
\\ (iii)  ${\rm Col}(S, P,h,\alpha)\cong {\rm Col}(S, P',h',\alpha')$
 via $G_\varphi((a,u))=(a,\varphi(u))$.
 \\ (iv)
 ${\rm Cay}(S,P,h,\alpha)\cong {\rm Cay}(S, P',h',\alpha')$
 via $H_\varphi((a,u))=(a,\varphi(u))$.
\end{proposition}

\begin{proof}
(i) We have
$h'(\varphi(u), \varphi(v\times_\alpha w) )=h(u,v\times_\alpha w )= \alpha(u\wedge v\wedge w)=\alpha'(\varphi(u) \wedge \varphi(v)\wedge \varphi(w))=h'(\varphi(u), \varphi(v)\times_{\alpha'} \varphi(w)).$
Since $h$ is nondegenerate it follows that $\varphi(u\times_{\alpha}v)=u\times_{\alpha'}v$.
\\ (ii)  is trivial.
\\ (iii) It is now easy to see that for $G=G_\varphi$, we have
$$G((a,u)(b,v))=G(ab-\frac{1}{2}(h(u,v)+\overline{h(u,v)}), va+u b+u\times_{\alpha}v)$$
$$= (ab-\frac{1}{2}(h(u,v)+\overline{h(u,v)}), \varphi(v)\varphi(a)+\varphi(u) \varphi( b)+\varphi(u\times_{\alpha}v))$$
and
$$G((a,u))G((b,v))$$
$$=(a,\varphi(u))(b,\varphi(v))=(ab-\frac{1}{2}(h'(\varphi(u),\varphi(v))+\overline{h'(\varphi(u),\varphi(v))}), \varphi(v)a+ \varphi(u) b+u\times_{\alpha'} v)$$
$$=(ab-\frac{1}{2}(h(u,v)+\overline{ h(u,v)}),\varphi(v))), \varphi(v)a+ \varphi(u) b+u\times_{\alpha'}v).$$
The second entries are equal due to our assumption that $\varphi$ is $S$-linear.
\\ (iv)
Analogously, we obtain  for $H=H_\varphi$:
$$H((a,u)(b,v))=H(ab-h(u,v), va+u \bar b+u\times_{\alpha}v)= (ab-h(u,v), \varphi(v)\varphi(a)+\varphi(u) \varphi( \bar b)+\varphi(u\times_{\alpha}v))$$
and
$$H((a,u))H((b,v))=(a,\varphi(u))(b,\varphi(v))=(ab-h'(\varphi(u),\varphi(v)), \varphi(v)a+ \varphi(u) \bar b+u\times_{\alpha'} v)$$
$$=(ab- h(u,v),\varphi(v))), \varphi(v)a+ \varphi(u) \bar b+u\times_{\alpha'}v).$$
\end{proof}

\end{document}
